\chapter{Σύστημα Οπτικοποίησης Αποτελεσμάτων}
\label{ch:dashboard}

Για την αποτελεσματική παρουσίαση και εξερεύνηση των αποτελεσμάτων αναπτύχθηκε
ένα διαδραστικό σύστημα οπτικοποίησης (\textit{dashboard}) με βάση
HTML5/JavaScript/Chart.js. Το σύστημα λειτουργεί ως αυτόνομο αρχείο, χωρίς
απαίτηση web server ή προγραμματιστικών εργαλείων, και είναι άμεσα
αναπαραγώγιμο.

% ─────────────────────────────────────────────────────────────────────────────
\section{Αρχιτεκτονική}
\label{sec:dash_arch}
% ─────────────────────────────────────────────────────────────────────────────

\subsection*{Τεχνολογίες}

\begin{itemize}
    \item \textbf{HTML5 + Vanilla JavaScript}: Το κύριο αρχείο \texttt{dashboard.html}
    (~2500 γραμμές κώδικα) δεν απαιτεί npm, webpack ή άλλα build tools.
    \item \textbf{Chart.js 4.x}: Βιβλιοθήκη JavaScript για διαδραστικά
    γραμμικά διαγράμματα, φορτωμένη μέσω CDN.
    \item \textbf{JSON data files}: 14 αρχεία δεδομένων (ένα ανά
    task × strategy × horizon) φορτώνονται δυναμικά μέσω \texttt{fetch()}.
    \item \textbf{Python http.server}: Απαιτείται μόνο για να εξυπηρετηθούν
    τα JSON αρχεία λόγω CORS policy· εκκινείται με
    \texttt{python -m http.server 8000}.
\end{itemize}

\subsection*{Δομή Δεδομένων (JSON Schema)}

Κάθε JSON αρχείο ακολουθεί ενιαίο σχήμα:
\begin{lstlisting}[language={}]
{
  "strategy": "cl" | "ol" | "mimo" | "direct",
  "task":     "price" | "load",
  "dates":    ["2025-12-01 00:00", ...],   // 744 timestamps
  "actual":   [110.2, 98.7, ...],          // πραγματικές τιμές
  "series": {
    "LightGBM":       [...],
    "XGBoost":        [...],
    "Ensemble-Best3": [...],
    "Naive-1":        [...],
    ...
  },
  "metrics": [
    {"Model": "LightGBM", "Type": "ml",
     "MAE": 10.334, "RMSE": 16.236, "sMAPE": 11.178},
    ...
  ]
}
\end{lstlisting}

% ─────────────────────────────────────────────────────────────────────────────
\section{Στοιχεία Διεπαφής Χρήστη}
\label{sec:dash_ui}
% ─────────────────────────────────────────────────────────────────────────────

\subsection*{Επιλογή Task και Στρατηγικής}

\begin{itemize}
    \item \textbf{Task toggle [PRICE / LOAD]}: Εναλλαγή μεταξύ τιμής DAM
    (€/MWh) και ηλεκτρικού φορτίου (MW).
    \item \textbf{Strategy buttons [CL | OL | MIMO | Direct]}: Επιλογή
    στρατηγικής· κάθε κουμπί φορτώνει το αντίστοιχο JSON.
    \item \textbf{Horizon sub-toggle [H=24 | H=168]}: Εμφανίζεται για
    OL, MIMO και Direct· επιλέγει daily (24h) ή weekly (168h) chunks.
\end{itemize}

\subsection*{Φιλτράρισμα και Ζουμ Χρονικής Περιόδου}

\begin{itemize}
    \item \textbf{Date range inputs}: Τα πεδία «Από» / «Έως» επιτρέπουν
    ζουμ σε υπό-περίοδο (π.χ. μία εβδομάδα Χριστουγέννων).
    \item \textbf{Δυναμική επαναϋπολογιστική μετρικών}: Μετά την εφαρμογή
    του date range, υπολογίζονται εκ νέου MAE/RMSE/sMAPE \textit{μόνο για
    την επιλεγμένη περίοδο}· ο υπότιτλος αλλάζει σε
    «\texttt{Metrics για επιλεγμένη περίοδο: \{from\} $\to$ \{to\}}''.
\end{itemize}

\subsection*{Φιλτράρισμα ανά Algorithm Family}

Το κουμπί \textbf{[Best] Best} ενεργοποιεί τη λογική \textit{best-per-family}:
για κάθε αλγόριθμο (LGBM, XGB, RF, SVR, MLP) εμφανίζεται μόνο η καλύτερη
παραλλαγή για την τρέχουσα στρατηγική. Αυτό αποφεύγει τον «σωρό» πολλαπλών
γραμμών του ίδιου αλγόριθμου και εστιάζει στη σύγκριση \textit{μεταξύ}
αλγόριθμων.

\subsection*{Χρωματική Κωδικοποίηση}

Κάθε αλγόριθμος / τύπος μοντέλου έχει σταθερό χρώμα σε όλες τις
στρατηγικές:

\begin{table}[htbp]
\centering
\caption{Χρωματική κωδικοποίηση dashboard.}
\label{tab:dash_colors}
\begin{tabular}{lll}
\toprule
\textbf{Μοντέλο} & \textbf{Χρώμα} & \textbf{Hex} \\
\midrule
LightGBM         & Sky blue       & \texttt{\#38bdf8} \\
XGBoost          & Orange         & \texttt{\#f97316} \\
Random Forest    & Green          & \texttt{\#22c55e} \\
MLP              & Purple         & \texttt{\#a855f7} \\
SVR              & Red            & \texttt{\#ef4444} \\
Ensemble Best3   & Gold           & \texttt{\#eab308} \\
Ensemble All     & Amber          & \texttt{\#fbbf24} \\
Actual           & White/Light    & \texttt{\#e5e7eb} \\
Naive-1          & Gray dashed    & \texttt{\#9ca3af} \\
Naive-24         & Gray dashed    & \texttt{\#6b7280} \\
Seasonal Profile & Pink dashed    & \texttt{\#ec4899} \\
\bottomrule
\end{tabular}
\end{table}

% ─────────────────────────────────────────────────────────────────────────────
\section{Χαρτογράφηση Αρχείων JSON}
\label{sec:dash_filemap}
% ─────────────────────────────────────────────────────────────────────────────

\begin{table}[htbp]
\centering\small
\caption{Χαρτογράφηση στρατηγικής/εργασίας σε αρχεία JSON.}
\label{tab:dash_filemap}
\begin{tabular}{lllp{5.5cm}}
\toprule
\textbf{Task} & \textbf{Strategy} & \textbf{Horizon} & \textbf{Αρχείο} \\
\midrule
price & cl & — & \texttt{...price\_cl\_monthly.json} \\
price & ol & h24 & \texttt{...price\_openloop\_h24\_monthly.json} \\
price & ol & h168 & \texttt{...price\_openloop\_h168\_monthly.json} \\
price & mimo & h24 & \texttt{...price\_mimo\_h24\_monthly.json} \\
price & direct & h24 & \texttt{...price\_direct\_h24\_monthly.json} \\
load  & cl & — & \texttt{...load\_cl\_monthly.json} \\
load  & ol & h24 & \texttt{...load\_openloop\_h24\_monthly.json} \\
\ldots & \ldots & \ldots & \ldots \\
\bottomrule
\end{tabular}
\end{table}

Συνολικά 14 JSON αρχεία (7 ανά task) φορτώνονται δυναμικά ανάλογα με την
επιλογή του χρήστη. Το μέγεθος κάθε αρχείου κυμαίνεται στα 200--500~KB.

% ─────────────────────────────────────────────────────────────────────────────
\section{Οδηγίες Χρήσης}
\label{sec:dash_usage}
% ─────────────────────────────────────────────────────────────────────────────

Για εκκίνηση του dashboard:
\begin{lstlisting}[language=bash]
# Από τον κατάλογο του project (epf_greece_starter/)
conda run -n epf python -m http.server 8000
# Άνοιγμα στον browser: http://localhost:8000/dashboard.html
\end{lstlisting}

\noindent Τυπική ροή χρήσης:
\begin{enumerate}
    \item Επιλέξτε εργασία (PRICE ή LOAD).
    \item Επιλέξτε στρατηγική (CL / OL / MIMO / Direct).
    \item Για OL/MIMO/Direct, επιλέξτε ορίζοντα (H=24 ή H=168).
    \item Πατήστε [Best] Best για εμφάνιση μόνο του καλύτερου ανά οικογένεια.
    \item Χρησιμοποιήστε Date Range για εστίαση σε υπό-περίοδο.
    \item Ο πίνακας μετρικών ανανεώνεται αυτόματα.
\end{enumerate}
